\subsection{Excluding order $58$}\label{sec:d31a58}
\begin{proposition}\label{prop:decic58-exclusion}
An elliptic curve over a number field of degree $10$ has no point of order $58$.
\end{proposition}
\begin{proof}
Put $Y=X_1(29)$ and $J=J_1(29)$, and let
$\pi:X_1(58)\longrightarrow Y$ forget the point of order $2$.
The curve $Y$ has genus $22$ and rational gonality $11$ \cite{DvH}.
By \cite[Theorems~3.1 and~4.13, Table~2]{DEHMZ},
\[
 J(\mathbb Q)\simeq C_4^2\times C_{64427244},
\]
and this group is generated by rational cuspidal classes.
Suppose $D$ is a noncuspidal effective rational divisor of degree $10$ on $X_1(58)$,
and put $D_Y=\pi_*D$.

The cusp signatures on $Y$ over $\mathbb Q$ and $\mathbb F_3$ are both
$1^{14}14$. If $\overline D_Y$ were cuspidal, it would lift to a
rational cuspidal divisor $C$ of degree $10$.
The torsion class $[D_Y-C]$ has zero reduction at $3$, so it is zero.
Thus $D_Y\sim C$, contrary to $\operatorname{gon}_{\mathbb Q}(Y)=11$,
unless $D_Y=C$. The latter is impossible since $D_Y$ is noncuspidal.

The noncuspidal source counts at $3$ through degree $10$ are
\[
 (0,0,0,0,0,140,224,756,2380,5796).
\]
Hence the noncuspidal part of $\overline D$ is a single closed point $Q_e$,
with multiplicity $1$ and $6\le e\le10$.
Write $\pi(Q_e)=P_d$ and $e=fd$.
The fibre is the degree-$3$ scheme $E[2]\setminus\{0\}$, so $f\in\{1,2,3\}$.
A quadratic fibre in degrees at most $10$ would lie over a base point of degree
$3$, $4$ or $5$.
All such base points have irreducible $2$-division cubics; the exact census
contains respectively $14$, $28$ and $14$ points in these degrees.
Thus $f\ne2$.
The only possible cubic fibre is $e=9,d=3$.
Otherwise $f=1$ and the curve at $P_d$ has rational nonzero $2$-torsion.
Every remaining cuspidal complement on $Y$ has degree at most $4$
and is supported on rational cusps.

Fix a rational cusp $c_0$, and put
$G=J(\mathbb F_3)$ and $H=\operatorname{red}_3J(\mathbb Q)$.
The proven class presentation gives
\[
 G\simeq C_4^3\times C_8^2\times C_{386563464},
 \qquad [G:H]=1536.
\]
Its factor basis consists of all $3535$ places of degree at most $9$,
which attain the unconditional generation bound.
The complete primary characters have image orders $32768$ and $9$,
and the rational cusps have image orders $64$ and $3$.
They therefore give an exact membership test for $H$.
The necessary conditions are
\[
 [P_d-dc_0]\in H\quad\text{or}\quad[3P_3-9c_0]\in H.
\]
The unchanged tests in degrees at most $9$ give $2324$ negative results.
They remain valid here because the larger cuspidal complements still
have zero image in $G/H$.

We compute every degree-$10$ place on $Y_{\mathbb F_3}$.
There are $5971$ places and $430$ diamond orbits, including seven orbits of size $7$.
Exactly $4060$ places, in $293$ orbits, have rational nonzero $2$-torsion.
For each representative $P$ we retain a function $f$ and an effective
factor-basis divisor $R$ satisfying
\[
 \operatorname{div}(f)+(32+s)c_0-P=R,\qquad \deg R=22+s.
\]
These exact identities compute its primary characters.
Since $H$ is diamond-stable, one representative per orbit suffices.
Exactly two orbits survive, of sizes $14$ and $7$.
Independent finite-field arithmetic reconstructs every marked point,
all $430$ orbits, the $2$-division cubics and all primary sums.
The source count $5796$ is recovered by counting the rational roots
of these cubics, and agrees with the independent Frobenius calculation.

The two surviving orbits are the reductions of the two known degree-$10$
points on $X_1(29)$ in van Hoeij's list \cite{DvH}.
Their residue fields have defining polynomials
\[
\begin{split}
 f_a(x)={}&x^{10}-2x^9+2x^8-5x^7+7x^6-4x^5\\
          &{}+4x^4-8x^3+5x^2-1,\\
 f_b(x)={}&x^{10}-5x^9+15x^8-31x^7+47x^6-56x^5\\
          &{}+51x^4-37x^3+20x^2-7x+1.
\end{split}
\]
The exact Tate coefficients and points of order $29$ are retained
in the certificate and checked afresh.
Both polynomials are irreducible modulo $3$.
The fourteen diamond translates of their traces reduce respectively
to the complete surviving orbits of sizes $7$ and $14$.
At $29$, the respective simple roots $11$ and $5$ give good reductions
with exactly $29$ points.
Consequently neither curve has rational nonzero $2$-torsion over its degree-$10$ field.

For a surviving reduction, choose the corresponding diamond translate $B$
of the known global trace.
The rational torsion class $[D_Y-B]$ reduces to zero at $3$, so it is zero.
Since $\deg D_Y=10<\operatorname{gon}_{\mathbb Q}(Y)$, we obtain $D_Y=B$.
Thus the underlying curve and residue field are those of the corresponding
known example. Its absence of rational $2$-torsion prevents a lift to $X_1(58)$.
This is a contradiction.
\end{proof}
